% Implementation & Setup -- owner: Kaiyuan Xu.  % Implementation & Setup -- owner: Kaiyuan Xu.  % Implementation & Setup -- owner: Kaiyuan Xu.  % Implementation & Setup -- owner: Kaiyuan Xu.  \input{setup} from the main report.

\section{Implementation and Experimental Setup}
\label{sec:setup}

\paragraph{System.}
All methods run in one Python package organised as a seven-stage pipeline:
maps, task streams and start locations, the simulator, an independent
collision checker, metrics, a parallel and resumable batch runner, and
summaries. One command runs a TOML configuration end to end, and every result
row records the git commit it was produced with. The planner (TP, TPTS and
space-time A$^*$), the token and the simulator are shared by every method;
a method is only a cost model plugged into A$^*$ (M0--M5) and, for the
assignment ablation, a task selector plugged into TP's choice of task
($S_{ij}$). Trust in the results rests on checks that run on every pull
request (331 tests): each run is replayed by the collision checker and
rejected on any vertex or swap conflict; the simulator validates every
planner decision; with $\lambda=\mu=0$ M3--M5 must reproduce TP's trajectories
exactly, and with $\beta=\gamma=0$ so must $S_{ij}$, under both TP and TPTS.

\paragraph{Maps.}
Both maps are $21\times35$ grids with 152 candidate start cells in the side
columns. \texttt{warehouse\_small} is the small warehouse of Ma et
al.~\cite{ma2017lifelong}, rebuilt cell by cell from their figure; its 200 task
endpoints form whole rows next to the shelves, so agents can walk along them
and each aisle is effectively three cells wide. \texttt{narrow\_corridor} is our
stress test with the same footprint: it removes the middle cross-aisle, so
every horizontal aisle is a 21-cell single-lane corridor entered only at its
ends, and turns the endpoint rows into 110 pockets cut into the shelf faces.
TP meets about twice as many head-on encounters per task there (1.35 vs.\ 0.67
at 50 agents and 5 tasks per step, test seeds).

\paragraph{Tasks and start locations.}
A run has 500 tasks. Task $k$ is released at $\lfloor k/f\rfloor$ for task
frequency $f$; its pickup and delivery are two distinct task endpoints drawn
uniformly at random. The paper does not say where agents start when there are
fewer than 50, so start cells are drawn from the 152 candidates. Both draws use
the run's seed, so every method sees exactly the same task stream and starts
for a given seed (common random numbers), which is what makes paired
comparisons possible.

\paragraph{Metrics.}
The service time of a task runs from its release to its delivery; we report
its mean, 95th percentile and maximum per run. We also record makespan,
throughput, total path cost, the wait fraction (share of an agent's busy
timesteps spent waiting), the Gini coefficient of per-cell visits, planning
time per timestep, and head-on encounters: two agents in the same row or
column, one or two cells apart with the cell between them (if any) passable
and empty, that both just moved one cell toward each other, counted once per
encounter and divided by the number of tasks.

\paragraph{Protocol.}
Parameters are tuned only on seeds 100--104 (30 agents, 1 and 5 tasks per
step), and a test keeps these disjoint from the test seeds 0--9. Each tuning
rule was written down before its runs were looked at: M1 and M2 take the lowest
mean service time; M3--M5 take, among settings within 0.5\% of the best mean
service time, the one with the lowest wait fraction; $S_{ij}$ takes, within the
same 0.5\% band, the one with the lowest 95th-percentile service time. The
chosen parameters are then fixed. The main experiment runs every method on
both maps with 10, 30 and 50 agents, 0.2, 1 and 5 tasks per step and seeds
0--9: 90 runs per method and map. A heavier-load experiment adds 80 agents at 1
and 5 tasks per step (20 runs per method and map). Every run is paired with TP
on the same (agents, frequency, seed); we report the mean paired change in
percent and a two-sided Wilcoxon signed-rank test over all pairs of a map, and
call a change significant at $p<0.05$.
 from the main report.

\section{Implementation and Experimental Setup}
\label{sec:setup}

\paragraph{System.}
All methods run in one Python package organised as a seven-stage pipeline:
maps, task streams and start locations, the simulator, an independent
collision checker, metrics, a parallel and resumable batch runner, and
summaries. One command runs a TOML configuration end to end, and every result
row records the git commit it was produced with. The planner (TP, TPTS and
space-time A$^*$), the token and the simulator are shared by every method;
a method is only a cost model plugged into A$^*$ (M0--M5) and, for the
assignment ablation, a task selector plugged into TP's choice of task
($S_{ij}$). Trust in the results rests on checks that run on every pull
request (331 tests): each run is replayed by the collision checker and
rejected on any vertex or swap conflict; the simulator validates every
planner decision; with $\lambda=\mu=0$ M3--M5 must reproduce TP's trajectories
exactly, and with $\beta=\gamma=0$ so must $S_{ij}$, under both TP and TPTS.

\paragraph{Maps.}
Both maps are $21\times35$ grids with 152 candidate start cells in the side
columns. \texttt{warehouse\_small} is the small warehouse of Ma et
al.~\cite{ma2017lifelong}, rebuilt cell by cell from their figure; its 200 task
endpoints form whole rows next to the shelves, so agents can walk along them
and each aisle is effectively three cells wide. \texttt{narrow\_corridor} is our
stress test with the same footprint: it removes the middle cross-aisle, so
every horizontal aisle is a 21-cell single-lane corridor entered only at its
ends, and turns the endpoint rows into 110 pockets cut into the shelf faces.
TP meets about twice as many head-on encounters per task there (1.35 vs.\ 0.67
at 50 agents and 5 tasks per step, test seeds).

\paragraph{Tasks and start locations.}
A run has 500 tasks. Task $k$ is released at $\lfloor k/f\rfloor$ for task
frequency $f$; its pickup and delivery are two distinct task endpoints drawn
uniformly at random. The paper does not say where agents start when there are
fewer than 50, so start cells are drawn from the 152 candidates. Both draws use
the run's seed, so every method sees exactly the same task stream and starts
for a given seed (common random numbers), which is what makes paired
comparisons possible.

\paragraph{Metrics.}
The service time of a task runs from its release to its delivery; we report
its mean, 95th percentile and maximum per run. We also record makespan,
throughput, total path cost, the wait fraction (share of an agent's busy
timesteps spent waiting), the Gini coefficient of per-cell visits, planning
time per timestep, and head-on encounters: two agents in the same row or
column, one or two cells apart with the cell between them (if any) passable
and empty, that both just moved one cell toward each other, counted once per
encounter and divided by the number of tasks.

\paragraph{Protocol.}
Parameters are tuned only on seeds 100--104 (30 agents, 1 and 5 tasks per
step), and a test keeps these disjoint from the test seeds 0--9. Each tuning
rule was written down before its runs were looked at: M1 and M2 take the lowest
mean service time; M3--M5 take, among settings within 0.5\% of the best mean
service time, the one with the lowest wait fraction; $S_{ij}$ takes, within the
same 0.5\% band, the one with the lowest 95th-percentile service time. The
chosen parameters are then fixed. The main experiment runs every method on
both maps with 10, 30 and 50 agents, 0.2, 1 and 5 tasks per step and seeds
0--9: 90 runs per method and map. A heavier-load experiment adds 80 agents at 1
and 5 tasks per step (20 runs per method and map). Every run is paired with TP
on the same (agents, frequency, seed); we report the mean paired change in
percent and a two-sided Wilcoxon signed-rank test over all pairs of a map, and
call a change significant at $p<0.05$.
 from the main report.

\section{Implementation and Experimental Setup}
\label{sec:setup}

\paragraph{System.}
All methods run in one Python package organised as a seven-stage pipeline:
maps, task streams and start locations, the simulator, an independent
collision checker, metrics, a parallel and resumable batch runner, and
summaries. One command runs a TOML configuration end to end, and every result
row records the git commit it was produced with. The planner (TP, TPTS and
space-time A$^*$), the token and the simulator are shared by every method;
a method is only a cost model plugged into A$^*$ (M0--M5) and, for the
assignment ablation, a task selector plugged into TP's choice of task
($S_{ij}$). Trust in the results rests on checks that run on every pull
request (331 tests): each run is replayed by the collision checker and
rejected on any vertex or swap conflict; the simulator validates every
planner decision; with $\lambda=\mu=0$ M3--M5 must reproduce TP's trajectories
exactly, and with $\beta=\gamma=0$ so must $S_{ij}$, under both TP and TPTS.

\paragraph{Maps.}
Both maps are $21\times35$ grids with 152 candidate start cells in the side
columns. \texttt{warehouse\_small} is the small warehouse of Ma et
al.~\cite{ma2017lifelong}, rebuilt cell by cell from their figure; its 200 task
endpoints form whole rows next to the shelves, so agents can walk along them
and each aisle is effectively three cells wide. \texttt{narrow\_corridor} is our
stress test with the same footprint: it removes the middle cross-aisle, so
every horizontal aisle is a 21-cell single-lane corridor entered only at its
ends, and turns the endpoint rows into 110 pockets cut into the shelf faces.
TP meets about twice as many head-on encounters per task there (1.35 vs.\ 0.67
at 50 agents and 5 tasks per step, test seeds).

\paragraph{Tasks and start locations.}
A run has 500 tasks. Task $k$ is released at $\lfloor k/f\rfloor$ for task
frequency $f$; its pickup and delivery are two distinct task endpoints drawn
uniformly at random. The paper does not say where agents start when there are
fewer than 50, so start cells are drawn from the 152 candidates. Both draws use
the run's seed, so every method sees exactly the same task stream and starts
for a given seed (common random numbers), which is what makes paired
comparisons possible.

\paragraph{Metrics.}
The service time of a task runs from its release to its delivery; we report
its mean, 95th percentile and maximum per run. We also record makespan,
throughput, total path cost, the wait fraction (share of an agent's busy
timesteps spent waiting), the Gini coefficient of per-cell visits, planning
time per timestep, and head-on encounters: two agents in the same row or
column, one or two cells apart with the cell between them (if any) passable
and empty, that both just moved one cell toward each other, counted once per
encounter and divided by the number of tasks.

\paragraph{Protocol.}
Parameters are tuned only on seeds 100--104 (30 agents, 1 and 5 tasks per
step), and a test keeps these disjoint from the test seeds 0--9. Each tuning
rule was written down before its runs were looked at: M1 and M2 take the lowest
mean service time; M3--M5 take, among settings within 0.5\% of the best mean
service time, the one with the lowest wait fraction; $S_{ij}$ takes, within the
same 0.5\% band, the one with the lowest 95th-percentile service time. The
chosen parameters are then fixed. The main experiment runs every method on
both maps with 10, 30 and 50 agents, 0.2, 1 and 5 tasks per step and seeds
0--9: 90 runs per method and map. A heavier-load experiment adds 80 agents at 1
and 5 tasks per step (20 runs per method and map). Every run is paired with TP
on the same (agents, frequency, seed); we report the mean paired change in
percent and a two-sided Wilcoxon signed-rank test over all pairs of a map, and
call a change significant at $p<0.05$.
 from the main report.

\section{Implementation and Experimental Setup}
\label{sec:setup}

\paragraph{System.}
All methods run in one Python package organised as a seven-stage pipeline:
maps, task streams and start locations, the simulator, an independent
collision checker, metrics, a parallel and resumable batch runner, and
summaries. One command runs a TOML configuration end to end, and every result
row records the git commit it was produced with. The planner (TP, TPTS and
space-time A$^*$), the token and the simulator are shared by every method;
a method is only a cost model plugged into A$^*$ (M0--M5) and, for the
assignment ablation, a task selector plugged into TP's choice of task
($S_{ij}$). Trust in the results rests on checks that run on every pull
request (331 tests): each run is replayed by the collision checker and
rejected on any vertex or swap conflict; the simulator validates every
planner decision; with $\lambda=\mu=0$ M3--M5 must reproduce TP's trajectories
exactly, and with $\beta=\gamma=0$ so must $S_{ij}$, under both TP and TPTS.

\paragraph{Maps.}
Both maps are $21\times35$ grids with 152 candidate start cells in the side
columns. \texttt{warehouse\_small} is the small warehouse of Ma et
al.~\cite{ma2017lifelong}, rebuilt cell by cell from their figure; its 200 task
endpoints form whole rows next to the shelves, so agents can walk along them
and each aisle is effectively three cells wide. \texttt{narrow\_corridor} is our
stress test with the same footprint: it removes the middle cross-aisle, so
every horizontal aisle is a 21-cell single-lane corridor entered only at its
ends, and turns the endpoint rows into 110 pockets cut into the shelf faces.
TP meets about twice as many head-on encounters per task there (1.35 vs.\ 0.67
at 50 agents and 5 tasks per step, test seeds).

\paragraph{Tasks and start locations.}
A run has 500 tasks. Task $k$ is released at $\lfloor k/f\rfloor$ for task
frequency $f$; its pickup and delivery are two distinct task endpoints drawn
uniformly at random. The paper does not say where agents start when there are
fewer than 50, so start cells are drawn from the 152 candidates. Both draws use
the run's seed, so every method sees exactly the same task stream and starts
for a given seed (common random numbers), which is what makes paired
comparisons possible.

\paragraph{Metrics.}
The service time of a task runs from its release to its delivery; we report
its mean, 95th percentile and maximum per run. We also record makespan,
throughput, total path cost, the wait fraction (share of an agent's busy
timesteps spent waiting), the Gini coefficient of per-cell visits, planning
time per timestep, and head-on encounters: two agents in the same row or
column, one or two cells apart with the cell between them (if any) passable
and empty, that both just moved one cell toward each other, counted once per
encounter and divided by the number of tasks.

\paragraph{Protocol.}
Parameters are tuned only on seeds 100--104 (30 agents, 1 and 5 tasks per
step), and a test keeps these disjoint from the test seeds 0--9. Each tuning
rule was written down before its runs were looked at: M1 and M2 take the lowest
mean service time; M3--M5 take, among settings within 0.5\% of the best mean
service time, the one with the lowest wait fraction; $S_{ij}$ takes, within the
same 0.5\% band, the one with the lowest 95th-percentile service time. The
chosen parameters are then fixed. The main experiment runs every method on
both maps with 10, 30 and 50 agents, 0.2, 1 and 5 tasks per step and seeds
0--9: 90 runs per method and map. A heavier-load experiment adds 80 agents at 1
and 5 tasks per step (20 runs per method and map). Every run is paired with TP
on the same (agents, frequency, seed); we report the mean paired change in
percent and a two-sided Wilcoxon signed-rank test over all pairs of a map, and
call a change significant at $p<0.05$.
